\documentclass[12pt, a4paper]{article}

% --- Thư viện hỗ trợ tiếng Việt và định dạng trang ---
\usepackage[utf8]{vietnam}          % Hỗ trợ tiếng Việt
\usepackage[margin=2.5cm]{geometry} % Căn lề 2.5cm đều 4 cạnh
\usepackage{graphicx}               % Chèn hình ảnh
\usepackage{amsmath, amssymb}       % Công thức toán học
\usepackage{booktabs}               % Tạo bảng đẹp
\usepackage{hyperref}               % Tạo link liên kết và mục lục tự động

\hypersetup{
    colorlinks=true,
    linkcolor=blue,
    filecolor=magenta,      
    urlcolor=cyan,
}

% --- Thông tin báo cáo ---
\title{\textbf{BÁO CÁO ĐỀ TÀI}\\ \Large Môn học / Dự án: Đề tài 5}
\author{\textbf{Người thực hiện:} Nguyễn Văn A \\ \textbf{MSSV:} 12345678}
\date{\today}

% --- Bắt đầu nội dung tài liệu ---
\begin{document}

% 1. Tạo trang bìa
\maketitle
\thispagestyle{empty}
\newpage

% 2. Tạo mục lục tự động
\tableofcontents
\newpage

% 3. Nội dung chính
\section{Giới thiệu tổng quan}
Đây là phần giới thiệu tổng quan về đề tài. Bạn có thể viết các đoạn văn mô tả bài toán, mục tiêu và phạm vi nghiên cứu tại đây.

\section{Cơ sở lý thuyết và Phương pháp}
\subsection{Cơ sở lý thuyết}
LaTeX hỗ trợ viết công thức toán học rất chuyên nghiệp. Ví dụ công thức nằm cùng dòng $E = mc^2$ hoặc công thức hiển thị riêng:

\begin{equation}
f(x) = \frac{1}{\sigma \sqrt{2\pi}} e^{-\frac{1}{2}\left(\frac{x-\mu}{\sigma}\right)^2}
\end{equation}

\subsection{Phương pháp thực hiện}
Dưới đây là danh sách các bước thực hiện:
\begin{itemize}
    \item \textbf{Bước 1:} Thu thập dữ liệu.
    \item \textbf{Bước 2:} Xử lý và làm sạch dữ liệu.
    \item \textbf{Bước 3:} Huấn luyện mô hình và đánh giá kết quả.
\end{itemize}

\section{Kết quả thực nghiệm}
\subsection{Bảng biểu}
Dưới đây là bảng tổng hợp kết quả đánh giá:

\begin{table}[htbp]
    \centering
    \caption{Bảng so sánh kết quả thực nghiệm}
    \label{tab:ketqua}
    \begin{tabular}{|c|l|c|}
        \hline
        \textbf{STT} & \textbf{Tên chỉ số} & \textbf{Giá trị} \\ \hline
        1 & Accuracy (Độ chính xác) & 95.5\% \\ \hline
        2 & Processing Time (Thời gian xử lý) & 1.25s \\ \hline
    \end{tabular}
\end{table}

\subsection{Hình ảnh}
Để chèn hình ảnh vào báo cáo, bạn lưu ảnh vào cùng thư mục với file `.tex` và sử dụng cú pháp:

\begin{figure}[htbp]
    \centering
    % \includegraphics[width=0.7\textwidth]{ten_anh.png}
    \caption{Sơ đồ hệ thống}
    \label{fig:sodo}
\end{figure}

\section{Kết luận và Phân tích}
Tóm tắt các kết quả đã đạt được, những hạn chế còn tồn tại và hướng phát triển tiếp theo của đề tài.

\end{document}